\lampiransection{DOKUMENTASI KINERJA NASI GERILYA DAN KOMPETITOR PADA GRABFOOD}
\label{app:dokumentasi-indeks-kinerja}

Lampiran ini memuat dokumentasi kinerja untuk penelitian \textit{Analisis Strategi Pemasaran Rumah Makan Nasi Gerilya pada Platform GrabFood Menggunakan Matriks SWOT}. Data yang disajikan mencakup jumlah penjualan bulanan dalam satuan bungkus, indeks \textit{Average Order Value} (AOV), indeks kunjungan harian, dan indeks kompetitor pada platform GrabFood. Jumlah penjualan bulanan disajikan dalam satuan bungkus sesuai grafik penjualan per bulan, sedangkan indeks AOV, indeks kunjungan harian, dan indeks kompetitor disajikan sebagai data turunan untuk membaca arah performa usaha dan posisi pembanding \textit{merchant}.

\begin{table}[!htbp]
    \centering
    \caption{Jumlah Penjualan Bulanan Nasi Gerilya}
    \label{tab:jumlah_penjualan_bungkus_grabfood}
    \begin{tabular}{|l|c|c|}
        \toprule
        \textbf{Periode} & \textbf{Jumlah Penjualan (Bungkus)} & \textbf{Perubahan} \\ \hline
        \midrule
        Des-24           & 6.702                               & --                 \\ \hline
        Jan-25           & 8.740                               & +30,4\%            \\ \hline
        Feb-25           & 8.911                               & +2,0\%             \\ \hline
        Mar-25           & 16.590                              & +86,2\%            \\ \hline
        Apr-25           & 15.129                              & -8,8\%             \\ \hline
        Mei-25           & 14.580                              & -3,6\%             \\ \hline
        Jun-25           & 13.467                              & -7,6\%             \\ \hline
        Jul-25           & 14.633                              & +8,7\%             \\ \hline
        Ags-25           & 14.259                              & -2,6\%             \\ \hline
        Sep-25           & 14.034                              & -1,6\%             \\ \hline
        Okt-25           & 13.009                              & -7,3\%             \\ \hline
        Nov-25           & 12.900                              & -0,8\%             \\ \hline
        \bottomrule
    \end{tabular}
    \tablesource{Dokumentasi grafik penjualan bulanan Nasi Gerilya; diolah peneliti (2026).}
\end{table}

\begin{table}[!t]
    \centering
    \caption{Indeks \textit{Average Order Value} (AOV)}
    \label{tab:indeks_aov_grabfood}
    \begin{tabular}{|l|c|c|}
        \toprule
        \textbf{Periode} & \textbf{Indeks} & \textbf{Perubahan} \\ \hline
        \midrule
        Jul-25           & 100,0           & --                 \\ \hline
        Ags-25           & 102,0           & +2,0\%             \\ \hline
        Sep-25           & 107,9           & +5,8\%             \\ \hline
        Okt-25           & 110,0           & +1,9\%             \\ \hline
        Nov-25           & 114,5           & +4,1\%             \\ \hline
        \bottomrule
    \end{tabular}
    \tablesource{Dokumentasi internal Nasi Gerilya; diolah peneliti (2026).}
\end{table}

\begin{table}[!t]
    \centering
    \caption{Indeks Kunjungan Harian}
    \label{tab:indeks_kunjungan_grabfood}
    \begin{tabular}{|l|c|c|}
        \toprule
        \textbf{Periode} & \textbf{Indeks} & \textbf{Perubahan} \\ \hline
        \midrule
        Sep-25           & 100,0           & --                 \\ \hline
        Okt-25           & 116,5           & +16,5\%            \\ \hline
        Nov-25           & 130,7           & +12,2\%            \\ \hline
        \bottomrule
    \end{tabular}
    \tablesource{Dokumentasi internal Nasi Gerilya; diolah peneliti (2026).}
\end{table}

\begin{table}[!t]
    \centering
    \small
    \caption{Indeks Kompetitor Nasi Padang Juli--Oktober 2025}
    \label{tab:indeks_kompetitor_grabfood}
    \begin{tabularx}{\textwidth}{|Y|>{\centering\arraybackslash}p{1.15cm}|>{\centering\arraybackslash}p{1.15cm}|>{\centering\arraybackslash}p{1.15cm}|>{\centering\arraybackslash}p{1.15cm}|>{\centering\arraybackslash}p{1.75cm}|}
        \toprule
        \textbf{\textit{Merchant}} & \textbf{Jul} & \textbf{Ags} & \textbf{Sep} & \textbf{Okt} & \textbf{\shortstack{$\Delta$ Jul--Okt}} \\ \hline
        \midrule
        Koto Gadang                & 100,0        & 105,1        & 86,2         & 94,2         & -5,8\%                                  \\ \hline
        Pondok Krakatau            & 100,0        & 110,1        & 89,0         & 107,5        & +7,5\%                                  \\ \hline
        Uda Sayang                 & 100,0        & 94,1         & 89,0         & 88,9         & -11,1\%                                 \\ \hline
        Dendeng Batokok            & 100,0        & 104,1        & 85,2         & 95,1         & -4,9\%                                  \\ \hline
        Istana Krakatau            & 100,0        & 119,2        & 108,4        & 119,3        & +19,3\%                                 \\ \hline
        Sederhana                  & 100,0        & 97,5         & 94,9         & 98,6         & -1,4\%                                  \\ \hline
        Gumarang Jaya              & 100,0        & 116,1        & 102,6        & 114,1        & +14,1\%                                 \\ \hline
        Nasi Gerilya               & 100,0        & 92,6         & 87,8         & 70,4         & -29,6\%                                 \\ \hline
        Garuda                     & 100,0        & 95,5         & 90,1         & 93,9         & -6,1\%                                  \\ \hline
        Pondok Gurih               & 100,0        & 102,9        & 97,0         & 101,5        & +1,5\%                                  \\ \hline
        \bottomrule
    \end{tabularx}
    \tablesource{Dokumentasi internal dan pembanding \textit{merchant}; diolah peneliti (2026).}
\end{table}

\FloatBarrier
Data pada Tabel~\ref{tab:jumlah_penjualan_bungkus_grabfood}, \ref{tab:indeks_aov_grabfood}, \ref{tab:indeks_kunjungan_grabfood}, dan \ref{tab:indeks_kompetitor_grabfood} digunakan sebagai dasar pembacaan arah performa usaha dan posisi pembanding pada platform GrabFood.
